% ============================================================
% DRAFT NOTE (discussion)
% Contains: Interpretation of where forecasting works, settled limitations,
% and a short conclusion.
% Written now: The limitations that are settled by design; interpretation
% awaits results.
% Pending: All result-dependent interpretation.
% Sources: pivot discussion; docs/_why/3_measuring_quality.md (quality vs
% divergence rationale)
% ============================================================
\section{Discussion and Limitations}

\todo{STALE SECTION: written for the divergence objective; to be rewritten for the quality-risk target. Do not review yet.}
\label{sec:discussion}

\todo{Interpret the held-out and transfer results: which horizons, tasks,
compressors, and ratios carry signal; where it degrades; what the baselines
already capture.}

Several limitations are settled by design. The target measures damage in
distribution space; the link to task quality is validated empirically in
Section~\ref{sec:setup} rather than guaranteed, and a run can diverge in
distribution while still completing its task. The claim is scoped to the
observed configuration grid, with compression fixed and active over each
forecast window: unseen compressors, unseen ratios, mid-window switches, and
prospective activation from an uncompressed state are all outside the
established result. The oracle label requires a paired uncompressed
evaluation on the same realized prefix, so supervision is only available
where such traces have been recorded. Primary supervision comes from one
model family; the transfer analysis tests, but does not guarantee, generality
beyond it. Finally, HERALD forecasts damage; it does not decide
interventions. Acting on the forecast, by adjusting ratio, disabling
compression, or re-prefilling, is downstream serving policy and out of scope
here.

\subsection*{Conclusion}
\todo{Two- or three-sentence conclusion once results are in: the question,
the evidence, and the deployment claim.}
